\section{Resultados e Discussão}
\label{sec:resultados}
\label{sec:discussao}

\subsection{Composição da amostra e contrastes principais}

A Tabela~\ref{tab:h1-principais} mostra que parte relevante da mudança de H1a antecede a inclusão das perdas. Na soja, o gradiente preço×capacidade passa de 0,08484 em A original a 0,04645 em A comum; no milho, de 0,05188 a 0,03090. Ao acrescentar intensidade e interação, B apresenta valores próximos aos da referência comum. A redução, portanto, não pode ser atribuída integralmente a um ajuste por quebra de safra. Ela acompanha a seleção de municípios e anos com histórico produtivo suficiente, cuja diferença de conversão média já aparecia na análise descritiva.

% Gerado por etapa9/publicar_tabelas.py; fontes e células registradas na auditoria.
\begin{table}[H]
\centering
\normalsize
\setlength{\tabcolsep}{2pt}
\renewcommand{\arraystretch}{1}
\normalcaption\captionstyle{\centering}
\caption{Conversão para agricultura: capacidade e perdas anteriores}
\label{tab:h1-principais}
\begin{tabular}{p{5.2cm}*{3}{p{3.3cm}}}
\hline
Variável & A original & A comum & B: intensidade \\
\hline\multicolumn{4}{l}{\textbf{Soja}} \\
$pZ$ & 0,08484 & 0,04645 & 0,04043 \\
 & (0,03578) & (0,05183) & (0,05307) \\
$L$ & -- & -- & -0,13128 \\
 &  &  & (0,06165) \\
$pL$ & -- & -- & 0,42338 \\
 &  &  & (0,29654) \\
$Z\tau$ & 0,04050 & 0,03728 & 0,03610 \\
 & (0,01444) & (0,02365) & (0,02380) \\
$R$ & -0,01348 & -0,00643 & -0,00650 \\
 & (0,00626) & (0,00437) & (0,00437) \\
Pares município--ano & 2.767 & 1.459 & 1.459 \\
Municípios & 134 & 106 & 106 \\
$R^2$ total & 0,3828 & 0,5109 & 0,5119 \\
$R^2$ \textit{within} & 0,0198 & 0,0159 & 0,0180 \\
\hline\multicolumn{4}{l}{\textbf{Milho}} \\
$pZ$ & 0,05188 & 0,03090 & 0,03218 \\
 & (0,01787) & (0,01823) & (0,01820) \\
$L$ & -- & -- & 0,04491 \\
 &  &  & (0,03955) \\
$pL$ & -- & -- & -0,01315 \\
 &  &  & (0,11571) \\
$Z\tau$ & 0,05382 & 0,05062 & 0,05080 \\
 & (0,01517) & (0,01791) & (0,01781) \\
$R$ & -0,01516 & -0,00322 & -0,00310 \\
 & (0,00671) & (0,00342) & (0,00344) \\
Pares município--ano & 2.767 & 2.376 & 2.376 \\
Municípios & 134 & 134 & 134 \\
$R^2$ total & 0,3807 & 0,4213 & 0,4217 \\
$R^2$ \textit{within} & 0,0165 & 0,0184 & 0,0190 \\
\hline\end{tabular}
\par\smallskip
\begin{minipage}{\textwidth}\normalsize\centering
Fonte: elaboração própria com MapBiomas 10.1, IBGE/PAM, GAEZ v4, Banco Mundial, CPI-U e CHIRPS.
\par\medskip
\raggedright Legenda: $p$: log-preço real centrado; $Z$: capacidade potencial padronizada; $L$: intensidade da perda; $\tau$: tendência em décadas; $R$: anomalia de chuva em centenas de milímetros. Preço em $t-1$ e perda em $t-2$. Erros-padrão municipais entre parênteses. Todas as colunas incluem efeitos municipais e de bioma por ano. A original: amostra disponível; A comum: mesmas chaves de B; B: inclusão de $L$ e $pL$.
\end{minipage}
\end{table}


Os coeficientes estão na escala da equação, e os valores entre parênteses são erros-padrão com agrupamento municipal. Todas as colunas incluem os efeitos municipais e de bioma por ano, a tendência por capacidade e a anomalia de chuva. O $R^2$ total incorpora a contribuição dos efeitos fixos, enquanto o indicador \textit{within} se refere à variação após sua absorção. A comparação desses ajustes entre amostras diferentes não constitui medida de superioridade da explicação econômica. A leitura substantiva se concentra nos contrastes e em sua precisão sob as três inferências.

Em B, os gradientes de H1a continuam positivos nas duas culturas, mas os intervalos de 95\% abrangem zero. O modificador de H1b é positivo na soja e próximo de zero, com sinal negativo, no milho; seus intervalos também abrangem zero nos três procedimentos. A versão binária com limiar de 20\% não produz contraste preciso. A Tabela~\ref{tab:h1-magnitudes} apresenta os intervalos dos contrastes principais. Esses resultados preservam a direção prevista para a capacidade no núcleo, sem confirmar sua positividade de forma geral; para as perdas, não permitem estabelecer uma modificação estável nem concluir que ela inexista.


A associação isolada da perda ajuda a esclarecer a diferença entre nível e interação. Na soja, seu coeficiente é negativo e rejeita zero com agrupamento municipal, mas não com agrupamento duplo ou HAC. Como o preço foi centrado, esse coeficiente se refere à associação da perda no patamar médio do logaritmo do preço. Mesmo que fosse preciso nos três casos, ele não demonstraria restrição financeira nem substituiria o contraste preço×perda. A mudança de precisão reforça a necessidade de distinguir o parâmetro examinado da narrativa que poderia explicá-lo.

% Gerado por etapa9/publicar_tabelas.py; fontes e células registradas na auditoria.
\begin{table}[H]
\centering
\normalsize
\setlength{\tabcolsep}{2pt}
\renewcommand{\arraystretch}{1}
\normalcaption\captionstyle{\centering}
\caption{Contrastes de capacidade e perda para aumento de 10\% no preço (p.p.)}
\label{tab:h1-magnitudes}
\begin{tabular}{p{2.4cm}*{4}{p{3.1cm}}}
\hline
Cultura / hipótese & Estimativa & IC95 mun. & IC95 mun.+ano & IC95 HAC(2) \\
Soja / H1a & 0,00385 & [-0,00618; 0,01388] & [-0,00666; 0,01437] & [-0,00273; 0,01044] \\
Soja / H1b & 0,00404 & [-0,00157; 0,00964] & [-0,00452; 0,01259] & [-0,00517; 0,01324] \\
Milho / H1a & 0,00307 & [-0,00036; 0,00650] & [-0,00139; 0,00752] & [-0,00075; 0,00688] \\
Milho / H1b & -0,00013 & [-0,00231; 0,00206] & [-0,00211; 0,00186] & [-0,00110; 0,00085] \\
\hline\end{tabular}
\par\smallskip
\begin{minipage}{\textwidth}\normalsize\centering
Fonte: elaboração própria com MapBiomas 10.1, IBGE/PAM, GAEZ v4, Banco Mundial, CPI-U e CHIRPS.
\par\medskip
\raggedright Legenda: H1a: um desvio-padrão adicional de capacidade; H1b: dez pontos percentuais adicionais de perda. p.p.: pontos percentuais. IC95: intervalo individual de 95\%. mun.: agrupamento municipal; mun.+ano: agrupamento duplo; HAC(2): covariância robusta à heterocedasticidade e à autocorrelação, banda dois.
\end{minipage}
\end{table}



A Tabela~\ref{tab:h1-magnitudes} traduz os coeficientes para uma alta de 10\% no preço, combinada a um desvio-padrão adicional de capacidade ou a dez pontos percentuais adicionais de perda. Para H1b, os contrastes são 0,00404 p.p. na soja e $-0,00013$ p.p. no milho. Com agrupamento duplo, os intervalos são $[-0,00452;\,0,01259]$ e $[-0,00211;\,0,00186]$ p.p. A unidade refere-se ao estoque nativo inicial, e não à variação percentual da área agrícola. As taxas médias de conversão nas amostras de B são 0,10810 e 0,08506 p.p.; elas contextualizam a escala, mas a divisão de um contraste por essas médias não o transforma em efeito causal proporcional.

A direção positiva dos gradientes de capacidade em B é compatível com a ideia de que uma valorização da cultura aumenta mais a receita potencial por hectare em locais produtivamente favorecidos. A imprecisão, entretanto, impede estabelecer essa relação de forma geral no recorte com histórico. A comparação com \citeonline{BermanEtAl2023} deve considerar que seu desenho, território e suporte de variação diferem dos utilizados aqui. Acrescentar perdas a um painel municipal não constitui, por si só, demonstração de um mecanismo ausente naquele estudo. O avanço específico está na distinção entre capacidade potencial, desempenho anterior e destino da conversão, acompanhada da verificação de quanto o resultado depende da composição e da inferência.

Essa distinção também limita a interpretação tecnológica. O potencial sob alto insumo e sequeiro representa uma possibilidade padronizada; não informa o insumo utilizado nas parcelas convertidas. Nos modelos antecedentes para soja, descritos na metodologia e estimados sem perdas, a associação positiva do cenário de alto insumo não se reproduz sob baixo insumo e se torna negativa com o índice de adequação. A ocupação agrícola anterior apresenta relação distinta do potencial. Esses resultados qualificam a leitura de H1a, em consonância com a heterogeneidade das respostas à intensificação discutida por \citeonline{Spera2017}. Como não incluem o contraste com perdas, permanecem antecedentes de H1a e não evidência de H1b.

\subsection{Estabilidade temporal e informação disponível}

As alternativas não revelam uma definição de perda que sustente H1b uniformemente. Nenhum dos dezesseis modificadores por cultura rejeita zero após o ajuste de Holm, em qualquer das três inferências. Há resultados pontualmente mais precisos: na soja, a janela desde 2006 produz modificador negativo e rejeição bruta no agrupamento duplo e no HAC. Após o ajuste na família, os p-valores passam a 0,5557 e 0,4161. O sinal oposto ao núcleo permanece pertinente, sem orientar a seleção do resultado principal.

A decomposição por períodos explicita essa instabilidade. Na soja, preço×perda é positivo em 2001--2005 e negativo nos dois períodos seguintes; no milho, é negativo no primeiro e pequeno nos demais. Os testes de nulidade conjunta e igualdade dos modificadores rejeitam com HAC após Holm nas duas culturas, sem uniformidade entre os procedimentos. Para H1a, o gradiente da soja em 2012--2024 é negativo, com intervalo de agrupamento duplo inteiramente abaixo de zero. Esse resultado contraria a direção da sub-hipótese naquele recorte e deve acompanhar a interpretação dos sinais positivos do núcleo. Nesse período, o coeficiente é $-0,04995$, com intervalo de $[-0,08465;\,-0,01524]$.

A informação desses contrastes está distribuída de forma desigual. Os períodos iniciais contêm cinco e seis preços anuais, e a soja reúne apenas cinco ocorrências de perda de pelo menos 20\% no primeiro. Na inferência HAC, mudar a banda entre zero e quatro mantém as rejeições brutas dos testes da soja; no milho, elas dependem da banda. Em B, Jangada em 2013, associada à safra de 2011, concentra 29,85\% da soma de quadrados residualizada de preço×perda da soja. Excluir esse par eleva o modificador, ainda impreciso; retirar 2003 muda seu sinal. A influência é documentada sem eliminar observações para obter uma direção desejada.

As verificações adicionais mantêm a incerteza. A reamostragem restrita por ano, com 999 replicações, produz p-valores de 0,603 e 0,877 para os modificadores de soja e milho. Nas 499 perturbações da referência histórica por cultura, os quantis de 2,5\% e 97,5\% também atravessam zero. São sensibilidades condicionais à construção do indicador, não intervalos que resolvem toda a incerteza. Na alternativa que substitui o preço defasado pelo futuro, o gradiente preço×capacidade do milho é positivo nas três inferências. Esse comportamento, limitado ao período com preço futuro disponível, recomenda cautela diante de tendências persistentes; não demonstra antecipação ou expectativa dos produtores.

A formulação sem sinal prévio de H1b encontra respaldo na diversidade de respostas possíveis, e não na equivalência de todas as explicações. Uma quebra pode reduzir recursos próprios, alterar acesso ao crédito e adiar investimento; também pode ser seguida de expansão se houver condições para financiar a atividade ou reorganizar a produção. O painel não observa essas decisões nem identifica qual restrição foi determinante. Por isso, o sinal positivo da interação na soja não prova recuperação de renda por abertura, assim como o sinal negativo em recortes posteriores não demonstra falta de liquidez. Essas interpretações permanecem hipóteses sobre processos que poderiam produzir as associações.

Essa leitura também delimita a comparação com \citeonline{OrdonezEtAl2025} e \citeonline{HsiaoEtAl2026}. Seus resultados de expansão após choques adversos não implicam que a interação entre uma cotação internacional e uma perda anterior deva ser positiva em Mato Grosso. De modo semelhante, as associações entre lucros setoriais e perda florestal encontradas por \citeonline{RichardsArima2026} não são diretamente comparáveis ao modificador de uma perda física municipal. O indicador deste artigo não observa o saldo financeiro da safra. A contribuição do contraste é mostrar o alcance e a precisão dessa relação específica, preservando como questões abertas a transmissão aos recursos e a escolha entre alternativas de uso da terra.

Para integrar risco agrícola e conservação, uma extensão futura precisaria observar a transmissão entre perda, recursos e decisões sobre a terra. Dados de seguro, crédito contratado, custos e receitas poderiam distinguir exposição produtiva de capacidade financeira, desde que sua posição temporal fosse explícita. Incluir renda ou crédito realizados depois da perda como controles pode retirar justamente parte do caminho que se deseja estudar. O desenho atual fornece critérios para organizar essa investigação, sem apresentar como efeito de política uma associação que ainda depende de choques omitidos e seleção de amostra.

\subsection{Destinos e extensões do contraste}

A Tabela~\ref{tab:h1-destinos-robustez} compara agricultura e pastagem nas mesmas observações de cada cultura. Os gradientes preço×capacidade são maiores para pastagem. A diferença entre destinos é imprecisa na soja e rejeita zero nos p-valores brutos das três inferências no milho. Neste último caso, +10\% no preço e um desvio-padrão adicional de capacidade correspondem a uma diferença de $-0,04315$ p.p. entre os gradientes de agricultura e pastagem. Na família de quatro contrastes, contudo, o p-valor ajustado permanece abaixo de 5\% apenas com agrupamento municipal. As diferenças de preço×perda são imprecisas mesmo antes do ajuste.

% Gerado por etapa9/publicar_tabelas.py; fontes e células registradas na auditoria.
\begin{table}[H]
\centering
\normalsize
\setlength{\tabcolsep}{2pt}
\renewcommand{\arraystretch}{1}
\normalcaption\captionstyle{\centering}
\caption{Gradientes por destino da conversão e diferença entre destinos}
\label{tab:h1-destinos-robustez}
\begin{tabular}{p{3.4cm}p{2.0cm}p{3.9cm}p{2.0cm}p{3.9cm}}
\hline
Destino & $\beta$ & IC95 de $\beta$ & $\delta$ & IC95 de $\delta$ \\
\hline\multicolumn{5}{l}{\textbf{Soja}} \\
Agricultura & 0,04043 & [-0,0699; 0,1508] & 0,42338 & [-0,4743; 1,3211] \\
Pastagem & 0,21472 & [-0,0065; 0,4360] & 0,28650 & [-3,2205; 3,7935] \\
Agricultura--pastagem & -0,17429 & [-0,4297; 0,0811] & 0,13688 & [-4,0375; 4,3112] \\
\hline\multicolumn{5}{l}{\textbf{Milho}} \\
Agricultura & 0,03218 & [-0,0145; 0,0789] & -0,01315 & [-0,2218; 0,1955] \\
Pastagem & 0,48490 & [0,0647; 0,9051] & -1,18909 & [-3,6083; 1,2302] \\
Agricultura--pastagem & -0,45272 & [-0,8824; -0,0231] & 1,17594 & [-1,2361; 3,5879] \\
\hline\end{tabular}
\par\smallskip
\begin{minipage}{\textwidth}\normalsize\centering
Fonte: elaboração própria com MapBiomas 10.1, IBGE/PAM, GAEZ v4, Banco Mundial, CPI-U e CHIRPS.
\par\medskip
\raggedright Legenda: $\beta$: coeficiente de preço×capacidade; $\delta$: coeficiente de preço×perda. IC95: intervalo individual de 95\%, com agrupamento municipal e anual. Mesmas chaves de B por cultura. A diferença entre destinos é estimada diretamente, incorporando sua covariância.
\end{minipage}
\end{table}



Os intervalos dessa tabela utilizam agrupamento por município e ano. O potencial permanece o da cultura, inclusive na equação de pastagem. A comparação permite reconhecer que os dois destinos apresentam associações diferentes em algumas especificações, sem converter capacidade agrícola em capacidade pecuária. Também não acompanha o deslocamento de atividades entre parcelas ou municípios.

A interação tripla com capacidade não torna H1b mais conclusiva. Seus coeficientes são positivos na soja e negativos no milho, mas nenhum dos seis testes da família rejeita após Holm nas três inferências. Os contrastes condicionais nos percentis de capacidade têm intervalos de agrupamento duplo que abrangem zero. Os intervalos dos gradientes com preço relativo à carne também abrangem zero. As restrições de coeficientes opostos não são rejeitadas, o que não estabelece equivalência entre medidas. Na amostra compartilhada pelas culturas, a inclusão conjunta não produz rejeição dos testes conjuntos de H1a ou H1b.

A comparação entre destinos ajuda a compreender por que produção e conversão não caminham necessariamente juntas. A substituição de pastagens por lavouras pode ampliar a produção sem abertura imediata de vegetação, enquanto novas pastagens podem surgir em outras áreas. \citeonline{MacedoEtAl2012} documentaram a importância de terras previamente abertas para a expansão da soja na porção florestal de Mato Grosso; \citeonline{Braganca2018} mostrou que preços relativos podem se associar a reorganização de usos sob condições particulares. O maior gradiente da pastagem em parte das estimativas é compatível com a relevância de examinar esses destinos em conjunto, mas não identifica a sequência de realocação descrita por tais mecanismos. As interdependências discutidas por \citeonline{SimoesEtAl2020} e \citeonline{JungEtAl2022} exigem informações que ultrapassam o destino municipal imediato.

\subsection{Heterogeneidade em torno dos marcos institucionais}

Nas datas principais, nenhum dos oito contrastes institucionais rejeita após Holm com agrupamento municipal ou duplo. No HAC, a interação de preço, perda, fração amazônica e período desde 2007 para o milho é negativa e mantém rejeição ajustada, assim como a alternativa desde 2006. Para +10\% no preço, +0,10 na perda e +0,10 na fração amazônica, o primeiro contraste equivale a $-0,00891$ p.p. na diferença pós--pré. A rejeição restrita ao HAC é apresentada como parte da evidência, sem ser tomada como concordância entre procedimentos. Para a data principal, os p-valores ajustados são 1,00000, 0,14054 e 0,00212, respectivamente, com agrupamento municipal, agrupamento duplo e HAC.

A variação que sustenta esse contraste está concentrada. Depois dos demais termos, permanece apenas 0,453\% da soma de quadrados da interação do milho após os efeitos fixos; um município concentra 19,0\% da variação parcial, e 2005, 41,9\%. As dezesseis matrizes institucionais de agrupamento duplo e a do modelo conjunto exigiram correção de autovalores negativos. O procedimento conserva os coeficientes e viabiliza os intervalos, mas não amplia a informação disponível. Na comparação pela cobertura florestal inicial em torno de 2012, o contraste H1a da soja é negativo, com rejeições brutas que não persistem após Holm.

A divergência inferencial no contraste do milho ilustra esse problema. A rejeição ajustada no HAC deve constar da interpretação, mas diz respeito à associação de perdas e preços num território potencialmente alcançado pelo pacto da soja. Não constitui efeito direto da Moratória sobre a comercialização de milho. A concentração da variação residual e a necessidade de corrigir matrizes de covariância reforçam o cuidado com sua precisão. Da mesma forma, o contraste pela floresta inicial não informa o custo de adequação ao Código em cada imóvel. Seria inadequado utilizar essas estimativas para concluir que uma das políticas conteve ou ampliou a resposta às quebras.

Essa delimitação não contradiz a evidência de conservação identificada por outros desenhos. \citeonline{HeilmayrEtAl2020} avaliaram a Moratória em áreas aptas à soja, e \citeonline{HardingEtAl2021} investigaram o condicionamento da resposta aos preços sob instrumentos de conservação. A comparação exige reconhecer diferenças no tratamento, na unidade de exposição e no contrafactual. Um efeito médio de redução da abertura pode coexistir com mudanças na sensibilidade aos incentivos; o presente artigo não dispõe da identificação necessária para separar essas duas relações como efeitos causais das políticas. Sua contribuição institucional é explicitar o alcance das medidas territoriais e o suporte empírico das interações.

\subsection{Implicações para a gestão ambiental}

Para a gestão ambiental, a implicação é acompanhar o caminho da expansão, além do volume produzido. O crescimento da safra pode resultar de rendimento, frequência de cultivo, substituição de uso ou incorporação de novas áreas; cada margem envolve pressões distintas. Uma leitura apoiada apenas em produção ou área agrícola líquida pode reunir processos opostos. Os dados de transição bruta oferecem uma medida mais próxima da abertura observada, mas ainda precisam ser relacionados às parcelas, à sequência de usos e às regras aplicáveis para sustentar decisões operacionais. O artigo não estima um limiar de preço para acionar fiscalização nem classifica municípios por risco previsto.

A utilidade gerencial dos resultados está, portanto, na organização de dimensões que precisam ser observadas conjuntamente. Preços situam os incentivos; capacidade e desempenho descrevem possibilidades e experiências locais; transições brutas mostram o destino imediato da abertura; informações por imóvel e cadeia permitiriam relacioná-los às obrigações e aos compromissos comerciais. A integração dessas dimensões pode orientar o desenho de monitoramento e a formulação de perguntas verificáveis. As estimativas apresentadas não estabelecem a eficácia de uma combinação específica de crédito, seguro e fiscalização, nem sustentam recomendações uniformes para os três biomas.

Por fim, a comparabilidade deve acompanhar qualquer uso territorial da evidência. O Pantanal possui suporte muito reduzido no núcleo, enquanto os municípios de cultivo recente podem ser excluídos pela falta de histórico. A alteração da amostra ao longo dos anos exige separar mudança econômica de mudança na informação disponível. Para além de Mato Grosso, essa é uma contribuição metodológica pertinente a regiões em que expansão agrícola, risco produtivo e conservação coexistem: reconhecer quem e quais decisões entram no contraste antes de transferir seu resultado para uma política. A validade dessa transferência depende das instituições, das alternativas de uso e da qualidade das observações em cada contexto.
